\documentclass[10pt,a4paper]{amsart}
\usepackage[margin=19mm]{geometry}
\usepackage{amsmath,amssymb,mathtools}
\usepackage[scr=rsfs,cal=euler]{mathalfa}
\usepackage{xfrac}
\usepackage[unicode,hidelinks,bookmarks=false]{hyperref}
\newcommand{\ie}{i.e.\ }
\newcommand{\cf}{cf.\ }
\newcommand{\resp}{respectively\ }
\newcommand{\Subsection}{\subsection}
\providecommand{\nobreakdash}{\nobreak\hbox{-}}
\setlength{\emergencystretch}{2em}
\multlinegap0pt
\usepackage{amssymb}
\usepackage{xcolor}
\newtheorem{example}{Example}
\newtheorem{lemma}{Lemma}
\title{Upper bounds for graded Betti numbers of projective schemes in the first nontrivial strand}
\author{Doyoon Ha, Minjae Kwon, JeongDon Lee, Jinhyung Park}
\date{Source: arXiv:2507.17231v2 (CC BY 4.0)}
\begin{document}
\maketitle

\noindent\emph{Source and licence.} Upper bounds for graded Betti numbers of projective schemes in the first nontrivial strand. Version arXiv:2507.17231v2 (CC BY 4.0), distributed by arXiv under the Creative Commons Attribution 4.0 International licence, http://creativecommons.org/licenses/by/4.0/, as recorded in the arXiv licence metadata for this exact version and shown on the abstract page https://arxiv.org/abs/2507.17231v2. Full licence notice: \texttt{licenses/arxiv-2507.17231-CC-BY-4.0.txt}.\par\smallskip

\noindent\textit{Editorial scope.} Counted unit: the article's lemma lem:bs-bound (source0.tex lines 427--440). The article's complete original proof of it is delivered with it (source0.tex lines 442--468, one \texttt{proof} environment of 27 lines). No mathematics is written, repaired or re-derived: the statement and the proof are the article's own lines, in the article's own notation, and no retained line is substituted or re-flowed.  Retained uncounted as the source-local context the proof needs: lines 354--376.  Nothing was omitted from any retained span, so the package carries no omission record.  No retained line contains a citation, so the wrapper carries no bibliography and no bibliography entry.  No retained line contains a figure environment, an image-inclusion command or drawing code, so no image is delivered and the package declares no asset dependency.  Editorial formatting only, in addition: an AMS \texttt{amsart} preamble that restates the article's own declarations and macros used by the retained lines, namely the environments \texttt{example}, \texttt{lemma} on separate counters, together with the standard packages those lines need.  The retained environments print in this document as Example 1, Lemma 1 in order of appearance, whereas the article prints the same items with the numbering of its own full document.  The complete native source of the article as distributed in the arXiv e-print for this exact version is bundled unchanged as \texttt{sources/182063/source0.tex}.
\begin{example}\label{ex:deqpure}
Consider the degree sequence $d^{e,q}:=(0,q+1, q+2, \ldots, q+e-1, q+e)$ of length $e$. Then
$$
\kappa_{p,q}(d^{e,q}) = {p+q-1 \choose q} {e+q \choose p+q}~~\text{ for $1 \leq p \leq e$}
~\text{ and }~
e(d^{e,q}) = {e+q \choose q}.
$$
The Betti table $\pi(d^{e,q})$ is the following:
\begin{center}
\texttt{ \begin{tabular}{c|cccccc}
         & $0$ & $1$ & $2$ & $\cdots$ & $e-1$ & $e$ \\ \hline
    $0$    & $1$ & $-$ &  $-$ & $\cdots$ & $-$ &$-$ \\
    $1$   & $-$ & $-$ & $-$ & $\cdots$ & $-$ & $-$ \\
    $\vdots$ & $\vdots$ &  $\vdots$ &  $\vdots$ &  $\ddots$ &  $\vdots$ &  $\vdots$ \\
    $q-1$ & $-$ & $-$ & $-$ & $\cdots$ & $-$ & $-$ \\
    $q$ & $-$ & ${q \choose q}{e+q \choose 1+q}$ & ${q+1 \choose q} {e+q \choose 2+q}$  & $\cdots$ & ${e+q-2 \choose q} {e+q \choose e-1+q}$ & ${e+q-1 \choose q} {e+q \choose e+q}$ 
\end{tabular}}
\end{center}

\smallskip

\noindent This is the Betti table of the $q$-secant variety of a variety of minimal degree. 
\end{example}

\begin{lemma}\label{lem:bs-bound}
Let $d:=(0, d_1, d_2, \ldots, d_e)$ be a degree sequence of length $e$ with $d_1 \geq q+1$ Then
$$
\kappa_{p,q}(d) \leq \underbrace{{p+q-1 \choose q} {e+q \choose p+q}}_{=\kappa_{p,q}(d^{e,q})}~~\text{ for $1 \leq p \leq e$}~\text{ and }~
e(d) \geq \underbrace{{e+q \choose q}}_{=e(d^{e,q})}.
$$
Furthermore, the following are equivalent:
\begin{enumerate}
\item $\kappa_{p,q}(d) = {p+q-1 \choose q} {e+q \choose p+q}$ for all $p$ with $1 \leq p \leq e$.
\item $\kappa_{p,q}(d) = {p+q-1 \choose q} {e+q \choose p+q}$ for some $p$ with $1 \leq p \leq e$.
\item $e(d) = {e+q \choose q}$.
\item $d=d^{e,q}$.
\end{enumerate}
\end{lemma}

\begin{proof}
See Example \ref{ex:deqpure} for the values of $\kappa_{p,q}(d^{e,q})$ and $e(d^{e,q})$. From this, we also get $(4) \Longrightarrow (1)$.
Since $d_1 \geq q+1$, it follows that $d_k \geq q+k$ for all $k \geq 1$. 
Then
$$
e(d)=\frac{1}{e!} \prod_{k \geq 1} d_k \geq \frac{1}{e!} \prod_{k \geq 1} (q+k) = e(d^{e,q}), 
$$
and the equality holds if and only if $d=d^{e,q}$. In particular, $(3) \Longleftrightarrow (4)$. Now, notice that $\kappa_{p,q}(d) \neq 0$ if and only if $d_p=q+p$. Assume that $d_p=q+p$ for $1 \leq p \leq m$. 
When $1 \leq k < p \leq m$, we have
$$
\frac{d_k}{d_p-d_k} = \frac{q+k}{p-k} = \frac{q+k}{|k-p|}.
$$
When $1 \leq p < k \leq e$ and $p \leq m$, we have $d_p=q+p$ and $d_k=q+k+a$ for some integer $a \geq 0$ so that
$$
\frac{d_k}{d_k-d_p} = \frac{q+k+a}{(k-p)+a} \leq \frac{q+k}{k-p} = \frac{q+k}{|k-p|}
$$
and the equality holds if and only if $d_k=q+k$.
Thus
$$
\kappa_{p,q}(d)=\prod_{\substack{k \geq 1 \\ k \neq p}} \frac{d_k}{|d_k - d_p|} \leq \prod_{\substack{k \geq 1 \\ k \neq p}} \frac{q+k}{|k-p|} = \kappa_{p,q}(d^{e,q})~~\text{ for $1 \leq p \leq m$},
$$
and the equality holds for some $p$ with $1 \leq p \leq m$ if and only if $d=d^{e,q}$. This also proves $(2) \Longleftrightarrow (4)$. We have shown that 
$$
(2) \Longleftrightarrow (3) \Longleftrightarrow (4) \Longrightarrow (1).
$$
The implication $(1) \Longrightarrow (2)$ is trivial.
\end{proof}

\end{document}
